\documentclass[10pt,a4paper]{amsart}
\usepackage[margin=19mm]{geometry}
\usepackage{amsmath,amssymb,mathtools}
\usepackage[scr=rsfs,cal=euler]{mathalfa}
\usepackage{xfrac}
\usepackage[unicode,hidelinks,bookmarks=false]{hyperref}
\newcommand{\ie}{i.e.\ }
\newcommand{\cf}{cf.\ }
\newcommand{\resp}{respectively\ }
\newcommand{\Subsection}{\subsection}
\providecommand{\nobreakdash}{\nobreak\hbox{-}}
\setlength{\emergencystretch}{2em}
\multlinegap0pt
\usepackage{bm}
\usepackage{bbm}
\usepackage{dsfont}
\usepackage{xspace}
\usepackage{cancel}
\usepackage{mathtools}
\usepackage{amsthm}
\makeatletter
\@ifundefined{thm}{\newtheorem{thm}{Thm}}{}
\@ifundefined{lem}{\newtheorem{lem}[thm]{Lemma}}{}
\makeatother
\newcommand{\set}[1]{\left\{ {#1} \right\}}
\title[A two-sorted theory of nilpotent Lie algebras]{A two-sorted theory of nilpotent Lie algebras}
\author{Christian d'Elb\'ee, Isabel M\"uller, Nicholas Ramsey,, Daoud Siniora}
\date{Source: arXiv:2407.12452v2 (CC BY 4.0)}
\begin{document}
\maketitle

\noindent\emph{Source and licence.} A two-sorted theory of nilpotent Lie algebras. Version arXiv:2407.12452v2 (CC BY 4.0), distributed under the Creative Commons Attribution 4.0 International licence, as recorded in the arXiv licence metadata for this exact version and shown on the abstract page https://arxiv.org/abs/2407.12452v2. Full rights notice: licenses/arxiv-2407.12452-CC-BY-4.0.txt.\par\smallskip

\noindent\textit{Editorial scope.} Counted unit: the Lem whose source label is \texttt{None} in the article (source0.tex lines 770--772, source label \texttt{None}), delivered together with the article's own complete original proof of it (source0.tex lines 774--805, one \texttt{proof} environment of 32 lines). No mathematics is written, repaired or re-derived: statement and proof are the article's own text line for line. Required context retained uncounted: extension of scalars (lines 628--630). Every definition, display and label of those lines is retained unchanged. Omitted from this package and not redistributed: 1--627, 631--769, 773--773, 806--1294. Nothing inside a retained excerpt is omitted or altered: every retained body is delivered exactly as the article's lines are written, so no omission receipt of any kind accompanies this package. Citations: the retained excerpts carry no citation command at all, so no bibliography entry of the article is transcribed and no reference is redistributed. Reference adaptation: none was needed and none was made; every reference inside the delivered unit resolves inside it, so no unresolved reference or citation is produced. Printed slips kept literal: no printed word, symbol, hypothesis or number of the counted statement or of its proof was altered. Layout and class: the article's own class is \texttt{amsart} with packages including enumerate, url, caption, xy, tabularx, adjustbox, amssymb, latexsym, amsmath, amsthm, amsfonts, mathrsfs, mathtools, tikz; the wrapper is the lane's pinned \texttt{amsart} editorial wrapper at 10pt on A4, which restates the theorem-like environments the retained text needs and adds the article's own macro definitions that the retained text needs (\texttt{\textbackslash set}), with extra wrapper packages (bm, bbm, dsfont, xspace, cancel, mathtools). The delivered numbering is the unit's own. Assets: no retained span contains a figure environment, a graphics-inclusion command, a PDF-page inclusion, a table or a TikZ picture; no image file is copied into this package, the package declares no asset dependency and no image file is redistributed with it. Licence: version arXiv:2407.12452v2 of this article is distributed under the Creative Commons Attribution 4.0 International licence, as recorded in the arXiv licence metadata for this exact version and shown on the abstract page https://arxiv.org/abs/2407.12452v2. A full rights notice is delivered at licenses/arxiv-2407.12452-CC-BY-4.0.txt. Characters: every retained excerpt is pure ASCII, so the delivered engine, XeLaTeX with its default Latin Modern text font, reports no missing character.
\begin{lem} \label{extension of scalars}
    Suppose $M =(K,V) \subseteq N=(K',V')$ are models of $T_{0}$ (or models of $T_{0}^{+}$ for a given theory of fields $T^{\dagger}$).  Then, in any model of $T_0$ containing $N$, we have $\langle K',V\rangle \cong (K', K' \otimes_{K} V)$, where $K' \otimes_{K} V$ denotes the $K'$-vector space obtained from $V$ by extension of scalars. 
\end{lem}

\begin{lem}
    The theory $T$ has quantifier elimination and its completions are determined by specifying the characteristic of the field.  More generally, if $T^{\dagger}$ is a complete theory extending the theory of fields with quantifier elimination, then $T^{+}$ is complete and has quantifier-elimination.  
\end{lem}

\begin{proof}
We first argue for $T$.  We use the well-known criterion for QE that if $M$ and $N$ are two $\aleph_{0}$-saturated models of $T$ whose fields have the same characteristic, then the set of partial isomorphisms from a finitely generated substructure of $M$ to a finitely generated substructure of $N$ has the back-and-forth property. Non-emptiness is easy, since if $K_{0}$ is the prime subfield then the unique map $f:(K_{0},\set{0})\to (K_{0},\set{0})$ is a partial isomorphism. 

Now suppose $M_{0} \subseteq M$ and $N_{0} \subseteq N$ are finitely generated substructures and $f : M_{0} \to N_{0}$ is an isomorphism. Suppose $a \in M \setminus M_{0}$.  Let $M_{1} = \langle a M_{0} \rangle$. We have to consider three cases:

\textbf{Case 1}:  $a \in K(M)$.  

Because $K(N)$ is algebraically closed, we can find some $b$ such that the isomorphism from $K(M_{0})$ to $K(N_{0})$ given by $f$ extends to an isomorphism from $K(M_{1})$, i.e. the subfield of $K(M)$ generated by $K(M_{0})a$, to the subfield of $K(N)$ generated by $K(N_{0})b$ and mapping $a \mapsto b$.  Let $N_{1} = \langle N_{0},b \rangle$.  Then, by Lemma \ref{extension of scalars}, $(V(M_{1}), [\cdot,\cdot])$ and $(V(N_{1}), [\cdot,\cdot])$ are obtained by extension of scalars from $K(M_{0})$ to $K(M_{1})$ and from $K(N_{0})$ to $K(N_{1})$ respectively.  The isomorphism $K(M_{1})$ to $K(N_{1})$ extending $f$, then, induces an isomorphism from $V(M_{1})$ to $V(N_{1})$ respecting the Lie algebra structure, as desired. 
%Note:  we're using here that extending the field only results in extension by scalars; this probably should be a separate lemma just for completeness.

\textbf{Case 2}: $a \in V(M)$ and $a$ is in the $K(M)$-span of $V(M_{0})$.  

Choose a basis $\overline{v}$ for $V(M_{0})$.  If $a$ is in the $K(M)$-span of $\overline{v} = (v_{0}, \ldots, v_{n-1})$ then, since $a \not\in V(M_{0})$, it must be the case that 
$$
a = \sum_{i < n} \alpha_{i} v_{i}
$$
with not all $\alpha_{i} \in K(M_{0})$.  Applying Case 1 at most $n$ times, we may extend $f$ to an isomorphism $f': \langle M_{0}, \alpha_{<n} \rangle \to \langle N_{0}, \beta_{<n} \rangle$ where $f'(\alpha_{i}) = \beta_{i}$ for all $i < n$.  Note that then
$$
f'(a) = \sum_{i < n} \beta_{i}f(v_{i}).
$$
Because we have the coordinate functions in the language, we have $M_{1} = \langle M_{0}, \alpha_{<n}\rangle$ and, setting $b = f'(a)$, we also have $N_{1} = \langle N_{0},b \rangle = \langle N_{0}, \beta_{<i} \rangle$ and $f' :M_{1} \to N_{1}$ is the desired extension. 

\textbf{Case 3}:  $a \in V(M)$ and $a$ is not in the $K(M)$-span of $V(M_{0})$.

We may assume that $[a,u]$ is in the $K(M)$-span of $V(M_{0})$ for all $u \in V(M_{0})$ (i.e that the Lie algebra over $K(M)$ spanned by $V(M_{0})$ is an ideal of the Lie algebra over $K(M)$ spanned $V(M_{1})$) since the general case follows from this one by reverse induction on the maximum $1\leq i\leq c+1$ such that $a\in P_i(M)\setminus P_{i+1}(M)$. Let $\{v_1,\dots, v_n\}$ be a $K(M_0)$-basis of $M_0$ and let $M_{0}'$ be the substructure of $M$ generated by $M_{0}$ and $\{[a,v_{i}] : i < n\}$.  Notice that $\overline{v}$ is still a basis of $V(M_{0}')$ though the field may grow. 
 By at most $n$ applications of Case 2, there is a structure $N_{0} \subseteq N_{0}' \subseteq N$ and an isomorphism $f': M_{0}' \to N_{0}'$ extending $f$.  

 Let $\overline{\alpha}$ be the sequence of structure constants for the subalgebra of $M$ spanned by $\overline{v}$ and $a$. By construction, $\overline{\alpha}$ is contained in $K(M_{0}')$ and hence $f'(\overline{\alpha})$ is contained in $K(N_{0}')$.  Since $f'(\overline{\alpha})$ is a sequence of structure constants for an LLA extending $f'(\overline{v})$, the axioms of $T$ entail that there is some $b \in V(N)$ such that $f'(\overline{v})$ and $b$ span an LLA with structure constants $f'(\overline{\alpha})$.  Then the map extending $f'$ and mapping $a \mapsto b$ defines an isomorphism from $M_{1} = \langle M_{0}'a \rangle$ to $\langle N_{0}'b \rangle$, which gives the desired extension.  

 The proof in the case that we are considering $T^{+}$ is essentially identical, with minor modifications. %First, to argue that the set of isomorphisms between finitely generated substructures is non-empty, we have that, for vectors $v \in V(M)$ and $w \in V(N)$, $\langle v \rangle$ and $\langle w \rangle$ are isomorphic as Lie algebras over the substructure of $K$ generated by $\emptyset$, which might not be the prime subfield \red{again $\set{0}\to \set{0}$ suffices?}. 
 Instead of considering all partial isomorphisms, we only consider those partial isomorphisms in which the induced $L^\dagger$-isomorphism of fields $K(M_0)\to K(N_0)$ is $T^\dagger$-elementary.  Then, since $K(M)$ and $K(N)$ are $\aleph_{0}$-saturated as models of $T^{\dagger}$, this map may be extended to incorporate a new field element, which yields Case 1. Case 2 and 3 are identical.
\end{proof}

\end{document}
